\subsection{Обобщённая производная по Соболеву}

\begin{definition}\label{def:sobolev-deriv}
Пусть $f \in \Lloc(\Om)$. Функция $g \in \Lloc(\Om)$ называется \textbf{обобщённой
производной $f$ по Соболеву}, $g = D^\alpha f$, если
\[
  D^\alpha T_f = T_g ,
\]
то есть
\[
  \int\limits_\Om g\,\varphi\, dx = (-1)^{|\alpha|} \int\limits_\Om f\, D^\alpha \varphi\, dx
  \qquad \forall \varphi \in C_0^\infty(\Om) .
\]
\end{definition}

\begin{remark}
Производная по Соболеву существует \textbf{не всегда}: распределение $D^\alpha T_f$
определено для любой $f \in \Lloc(\Om)$, но может оказаться сингулярным — тогда
никакой функции $g$ ему не соответствует.
\end{remark}

\inputfig{figures/fig-heaviside}

\begin{example}[функция Хевисайда]\label{ex:heaviside}
Пусть $\Om = \R$ и
\[
  H(x) = \begin{cases} 0, & x < 0, \\ 1, & x \ge 0. \end{cases}
\]
Тогда
\begin{multline*}
  D T_H(\varphi) = -\,T_H(\varphi')
  = -\int\limits_{-\infty}^{+\infty} H(x)\,\varphi'(x)\, dx \\
  = -\int\limits_0^{+\infty} \varphi'(x)\, dx
  = -\big(0 - \varphi(0)\big) = \varphi(0) = \delta(\varphi) ,
\end{multline*}
то есть $H' = \delta$. Распределение $\delta$ сингулярно (пример~\ref{ex:delta}),
значит \textbf{у функции Хевисайда нет соболевской производной}.
\end{example}

\begin{example}\label{ex:x2sign}
Пусть
\[
  f(x) = x^2 \sgn x = \begin{cases} x^2, & x \ge 0, \\ -x^2, & x < 0. \end{cases}
\]
Тогда
\[
  f'(x) = \begin{cases} 2x, & x \ge 0, \\ -2x, & x < 0, \end{cases} = 2|x| ,
  \qquad
  f''(x) = \begin{cases} 2, & x \ge 0, \\ -2, & x < 0, \end{cases} = 2\sgn x ,
\]
и обе эти производные — обычные функции из $\Lloc(\R)$. А вот третья производная уже
сингулярна:
\[
  f'''(\varphi) = -\,T_{f''}(\varphi') = -\int\limits_{-\infty}^{+\infty} 2\sgn x \cdot \varphi'(x)\, dx
  = 2\int\limits_{-\infty}^{0}\varphi'\,dx - 2\int\limits_0^{+\infty}\varphi'\,dx
  = 4\varphi(0) ,
\]
то есть $f''' = 4\delta$.
\end{example}

\begin{idea}
Видна общая картина: каждое дифференцирование «съедает» одну степень гладкости, и как
только у функции появляется скачок, следующая производная выпадает из $\Lloc$ и становится
сингулярной. Функция из примера~\ref{ex:x2sign} лежит в $\Wp{p}{2}$, но не в $\Wp{p}{3}$.
\end{idea}

\subsubsection*{Свойства обобщённой производной}

\begin{proposition}[свойство 1: единственность]\label{prop:deriv-unique}
Обобщённая производная единственна с точностью до множества меры нуль: если $g_1$ и $g_2$
обе удовлетворяют определению~\ref{def:sobolev-deriv}, то $g_1 = g_2$ почти всюду в $\Om$.
\end{proposition}

\begin{proof}
Вычитая два равенства, получаем
\[
  \int\limits_\Om (g_1 - g_2)\,\varphi\, dx = 0 \qquad \forall \varphi \in C_0^\infty(\Om) ,
\]
откуда по лемме Дюбуа--Реймона (лемма~\ref{lem:dbr}) $g_1 - g_2 = 0$ п.\,в.
\end{proof}

\begin{proposition}[свойство 2: линейность]\label{prop:deriv-linear}
Пусть $f_1, f_2 \in \Lloc(\Om)$ и существуют $D^\alpha f_1, D^\alpha f_2 \in \Lloc(\Om)$.
Тогда для любых $a_1, a_2 \in \R$ существует
$D^\alpha(a_1 f_1 + a_2 f_2) \in \Lloc(\Om)$ и
\[
  D^\alpha(a_1 f_1 + a_2 f_2) = a_1 D^\alpha f_1 + a_2 D^\alpha f_2 .
\]
\end{proposition}

\begin{proof}
Складывая определяющие равенства с коэффициентами $a_1, a_2$, получаем для всех
$\varphi \in C_0^\infty(\Om)$
\[
  \int\limits_\Om \big(a_1 D^\alpha f_1 + a_2 D^\alpha f_2\big)\varphi\, dx
  = (-1)^{|\alpha|} \int\limits_\Om \big(a_1 f_1 + a_2 f_2\big) D^\alpha \varphi\, dx ,
\]
а это и есть определение~\ref{def:sobolev-deriv} для функции $a_1 f_1 + a_2 f_2$.
\end{proof}

\subsubsection*{Сходимость в $\Lloc(\Om)$}

\begin{definition}
$f_k \to f$ в $\Lloc(\Om)$, если
\[
  \int\limits_{\Om'} |f_k - f|\, dx \xrightarrow[k \to \infty]{} 0
  \qquad \forall\, \Om' \Subset \Om ,
\]
то есть сходимость в $\Lp{1}(\Om')$ на каждом компактно вложенном подмножестве.
\end{definition}

\begin{remark}[метрика]
Эта сходимость метризуема. Возьмём $\Om_k \Subset \Om$ с $\bigcup_k \Om_k = \Om$ и
положим
\[
  \rho_k(f_1, f_2) = \int\limits_{\Om_k} |f_1 - f_2|\, dx ,
  \qquad
  \rho(f_1, f_2) = \sum_{k=1}^{\infty} \frac{1}{2^k}\,
  \frac{\rho_k(f_1, f_2)}{1 + \rho_k(f_1, f_2)} .
\]
Ряд сходится, так как каждое слагаемое не превосходит $2^{-k}$, и сумма не больше $1$.
\end{remark}

\begin{idea}
Дробь $\frac{t}{1+t}$ обрезает каждое слагаемое до величины меньше единицы — без неё
ряд мог бы расходиться, ведь $\rho_k$ не ограничены. При этом $\frac{t}{1+t} \to 0$
тогда и только тогда, когда $t \to 0$, так что сходимость по метрике $\rho$ равносильна
сходимости всех $\rho_k$ к нулю, то есть сходимости на каждом $\Om' \Subset \Om$.
\end{idea}

\subsubsection*{Свойство 3: признак существования обобщённой производной}

\begin{theorem}[признак существования производной по Соболеву]\label{thm:crit-deriv}
Пусть
\[
  f_k \in \Lloc(\Om) \quad \forall k, \qquad
  \exists\, D^\alpha f_k \in \Lloc(\Om) \quad \forall k,
\]
\[
  f_k \to f \ \text{ в } \Lloc(\Om), \qquad
  D^\alpha f_k \to g \ \text{ в } \Lloc(\Om) .
\]
Тогда
\[
  \exists\, D^\alpha f \in \Lloc(\Om): \qquad D^\alpha f = g .
\]
\end{theorem}

\begin{proof}
Пусть $\varphi \in C_0^\infty(\Om)$. По определению обобщённой производной для каждого $k$
\[
  \int\limits_\Om f_k\, D^\alpha \varphi\, dx
  = (-1)^{|\alpha|} \int\limits_\Om D^\alpha f_k \cdot \varphi\, dx ,
\]
а так как $\varphi$ финитна, интегрирование идёт лишь по $\supp \varphi$:
\[
  \int\limits_{\supp \varphi} f_k\, D^\alpha \varphi\, dx
  = (-1)^{|\alpha|} \int\limits_{\supp \varphi} D^\alpha f_k \cdot \varphi\, dx .
\]

\emph{Предел в левой части.}
\begin{multline*}
  \left| \int\limits_{\supp \varphi} f_k\, D^\alpha \varphi\, dx
       - \int\limits_{\supp \varphi} f\, D^\alpha \varphi\, dx \right|
  \le \int\limits_{\supp \varphi} |f_k - f|\cdot |D^\alpha \varphi|\, dx \\
  \le \sup_{\R^n} |D^\alpha \varphi| \int\limits_{\supp \varphi} |f_k - f|\, dx
  \xrightarrow[k \to \infty]{} 0 ,
\end{multline*}
поскольку $f_k \to f$ в $\Lloc(\Om)$, а $\supp \varphi \Subset \Om$ — компакт.

\emph{Предел в правой части} получается аналогично:
\[
  \int\limits_{\supp \varphi} D^\alpha f_k \cdot \varphi\, dx
  \xrightarrow[k \to \infty]{}
  \int\limits_{\supp \varphi} g\, \varphi\, dx .
\]

Тогда перейдём к пределу:
\[
  \int\limits_{\supp \varphi} f\, D^\alpha \varphi\, dx
  = (-1)^{|\alpha|} \int\limits_{\supp \varphi} g\, \varphi\, dx
  \quad\Longrightarrow\quad
  \int\limits_\Om f\, D^\alpha \varphi\, dx
  = (-1)^{|\alpha|} \int\limits_\Om g\, \varphi\, dx ,
\]
то есть $g = D^\alpha f$ на $\Om$.
\end{proof}

\begin{idea}
Признак позволяет получать обобщённые производные предельным переходом: достаточно
подобрать удобную последовательность гладких (или просто более хороших) функций,
сходящуюся вместе с производными. Сходимость нужна лишь локально, в $\Lloc$, —
это самое слабое требование, какое ещё даёт результат.
\end{idea}

\begin{example}[$(|x|)' = \sgn x$]
Знаем, что $(|x|)' = \sgn x$ — проверим это признаком.
\end{example}

Возьмём $f = |x|$ и последовательность $f_k$, совпадающую с $|x|$ вне
$\left[-\tfrac1k, \tfrac1k\right]$ и сглаженную внутри этого отрезка.
Тогда $f_k \to f$, а производные $f_k'$ — кусочно-линейные функции, переходящие
от $-1$ к $1$ на $\left[-\tfrac1k, \tfrac1k\right]$, так что $f_k' \to \sgn x$.
По признаку~\ref{thm:crit-deriv} получаем то же самое: $(|x|)' = \sgn x$.

\inputfig{figures/fig-abs-x}
